\documentclass[conference]{IEEEtran}
\usepackage{cite}
\usepackage{amsmath,amssymb,amsfonts}
\usepackage{graphicx}
\usepackage{textcomp}
\usepackage{xcolor}
\usepackage{booktabs}

% TODO markers use this so they're visually obvious in draft PDFs
\newcommand{\todo}[1]{\textcolor{red}{\textbf{[TODO: #1]}}}

\begin{document}

\title{Payment-State-Aware Root Cause Analysis for Cascading Failures in
Event-Driven Payment Processing Pipelines}

\author{
\IEEEauthorblockN{Sarvadnya \todo{full author list, affiliations}}
\IEEEauthorblockA{\todo{institution}}
}

\maketitle

\begin{abstract}
We present ClearFlow, a live, running 8-service event-driven ISO~20022
payment pipeline instrumented so that transactional state -- liquidity
reservations, AML holds, idempotency-key collisions, and settlement
finality -- is queryable per payment, and a fault-injection harness that
triggers \emph{real} faults against this running system (process crashes,
real AML holds, real idempotency collisions) rather than synthesizing
labeled ground truth. Existing microservice root-cause-analysis (RCA)
methods reason over generic telemetry and topology and do not model this
payment-specific state as a causal factor. We evaluate a payment-state-aware
RCA method against a telemetry/topology baseline, a naive telemetry-only
baseline, and a general-purpose LLM given identical topology/telemetry
evidence, using AC@k with Wilson confidence intervals on a held-out set of
25 live incidents disjoint from the incidents used to develop the method.
The payment-state-aware method reaches AC@1 0.60 (95\% CI 0.41--0.77)
versus 0.48 and 0.44 for the topology and telemetry baselines respectively,
and is statistically significant against a random-guessing floor
($p=0.0034$); the LLM baseline, given equivalent evidence, scores 0.28 --
below both non-LLM baselines. An ablation isolates which mechanism drives
this result rather than reporting it as a black box: the six
payment-domain-specific state fractions (AML hold, liquidity dwell,
idempotency, settlement failure, validation stall) alone match the full
method exactly (0.60), while a generalized process-crash completion-stall
signal added to extend coverage to non-payment-domain faults contributes
nothing extra and, used alone, underperforms even the topology baseline
(0.40). The validated claim is accordingly narrower than the headline
number: payment-domain transactional state, not process-crash-derived
evidence, is what this evaluation supports as a causal RCA signal, and
this is also why the method's dev-set improvement on confounded incidents
does not replicate in the held-out set ($n=3$) -- a limitation reported
directly rather than omitted. Two implementation-level findings generalize
beyond this system regardless: a single-snapshot state flag is not
specific evidence of a stuck fault without a dwell-time gate, and a
process-crash fault produces no domain-state anomaly at all, requiring
evidence framed as ``did this payment reach the next stage's completion
event,'' not as a state value -- even though, on this evaluation's own
evidence, that reframing alone is not sufficient to recover accuracy on
crash-based faults. A direct evidence-vs-reasoning experiment gives a
clean answer to why the LLM-based methods trail the formula: a static
LLM given the identical G0--G3 payment-state evidence the formula uses
scores exactly the same as one given only G0--G2 evidence (0.28 vs.\
0.28, zero net change) -- richer evidence alone does not help a static
LLM. An agentic variant given real tool access to the live system
instead of a static summary reaches 0.44--0.52, more than double the
static baseline, from the same underlying information but a reasoning
mode able to actively decide what to inspect. Separately, an
interventional experiment no
passive-dataset methodology in this literature can run -- triggering the
identical fault under two different real, controlled payment-state
backgrounds -- finds no detectable difference in cascade propagation at
the sample size tested ($n=4$--$5$), reported as an honest null result
with its own limitations stated rather than omitted or oversold.
\end{abstract}

\begin{IEEEkeywords}
root cause analysis, cascading failures, event-driven architecture, payment systems, microservices, observability
\end{IEEEkeywords}

\section{Introduction}
\label{sec:intro}

Payment processing pipelines are distributed systems with a property most
microservice architectures do not have: every request carries
\emph{transactional state} that persists and constrains behavior long after
the request itself has been handled. A liquidity reservation taken against a
nostro account blocks other payments from that account until it is released.
An AML hold on a sanctioned-entity match stops a payment from progressing
regardless of whether every downstream service is healthy. Once a settlement
is finalized, it cannot be rolled back, only compensated. These are not
implementation details -- they are causal factors in how a failure in one
stage cascades to others, and they have no analogue in a generic web service
or content-delivery pipeline.

Existing cascading-failure root-cause-analysis (RCA) work treats
microservices generically. Soldani et al.~\cite{soldani2021} explain
cascading failures declaratively but reason over service dependencies, not
domain state. Recent 2025--2026 systems push RCA accuracy substantially:
DynaCausal~\cite{dynacausal2025} reaches an average AC@1 of 0.63 by modeling
time-varying spatio-temporal dependencies across logs, traces, and metrics;
KRCA~\cite{krca2026} reaches AC@1 0.88 in production at Kuaishou via a
multi-agent causal-graph pipeline. Neither models transactional/domain state
as a causal factor -- both reason over infrastructure-level signals
(latency, error rate, trace topology) that a payment-specific fault can leave
completely undisturbed. The closest prior work,
ServiceGraph-FM~\cite{servicegraphfm2026}, targets payment systems directly
and models anomaly propagation as diffusion over the payment-service
dependency graph with causal attention over multi-hop paths -- but this is
still a \emph{structural} model of which services call which; it does not
represent liquidity-lock, AML-hold, idempotency, or settlement-finality
state as first-class causal evidence. This paper's central claim is that
this gap is not cosmetic: Section~\ref{sec:eval} shows state-agnostic
methods, including a topology-aware heuristic and a general-purpose LLM
given the same telemetry evidence, are actively misled on exactly the
incident class where a louder downstream symptom masks the true upstream
root cause -- and that payment-state evidence recovers the correct
diagnosis where telemetry alone cannot.

\textbf{Contributions:}
\begin{itemize}
  \item A live, running 8-service ISO~20022 payment pipeline
  (Section~\ref{sec:system}) instrumented so that liquidity-lock, AML-hold,
  idempotency, and settlement-finality state is queryable per payment, not
  just logged as an implementation detail.
  \item A fault-injection harness that triggers \emph{real} faults against
  this live system -- process crashes, real AML holds, real idempotency
  collisions -- rather than synthesizing labeled ground truth, producing
  cascades whose propagation, severity, and affected-payment counts are
  measured post-hoc from real Elasticsearch/Kafka evidence, not asserted in
  advance (Section~\ref{sec:faults}).
  \item A payment-state-aware RCA method evaluated against a state-agnostic
  topology/telemetry baseline and a general-purpose LLM baseline given
  identical evidence, using AC@k with Wilson confidence intervals and a
  held-out validation split disjoint from the incidents used to develop the
  method (Section~\ref{sec:eval}) -- reporting the result honestly,
  including a self-audit of methodological weaknesses found and fixed
  during evaluation (Section~\ref{sec:discussion}).
\end{itemize}

\section{Related Work}
\label{sec:related}

\subsection{Cascading Failure RCA in Microservices}
Soldani et al.~\cite{soldani2021} explain cascading failures in
microservice applications declaratively from a static dependency model --
general-purpose, domain-agnostic, and, notably for the fault taxonomy in
Section~\ref{sec:faults}, not built to distinguish a genuine upstream root
cause from a louder downstream symptom under backpressure.
RCAEval~\cite{rcaeval2024} (735 failure cases, 3 systems, 11 fault types)
and its fault-propagation-aware successor~\cite{fpa2025} (1430 validated
failures, 25 fault types, hierarchical/impact-driven ground truth)
established AC@k/Avg@k as the standard evaluation protocol for this line of
work; Section~\ref{sec:eval} adopts AC@k directly rather than inventing a
new metric, while departing from both benchmarks' methodology in one
respect worth stating plainly: neither evaluates against a real, running
system triggering real faults -- both, like the vast majority of RCA
benchmarks including this paper's own earlier synthetic pass, generate
labeled ground truth rather than measuring it from a live incident.
Fang et al.~\cite{fpa2025} report a finding this paper's own evaluation
independently reproduces in a different form: on 4 widely-used public RCA
benchmarks, simple rule-based methods match or beat sophisticated SOTA
models, because benchmark and evaluation design can make sophistication
look like an advantage when it is not actually being tested by the
benchmark. Section~\ref{sec:eval}'s LLM baseline shows the same pattern --
a general-purpose language model given identical evidence to a formulaic
heuristic scores \emph{worse} than the heuristic, not better.
DynaCausal~\cite{dynacausal2025} and KRCA~\cite{krca2026} represent the
current state of the art on generic microservice telemetry (average AC@1
0.63 and 0.88 respectively, the latter in production), and both attribute
their gains to modeling \emph{temporal} and \emph{structural} dynamics more
faithfully -- richer signal within the same evidence class this paper's
\texttt{loudest\_metric\_baseline} and \texttt{graph\_topology\_baseline}
occupy (Section~\ref{sec:method}), not a different evidence class. None of
this line of work represents transactional/domain state as a causal factor.

\textbf{Closest prior art.} ServiceGraph-FM~\cite{servicegraphfm2026}
targets payment-system RCA directly: a pretrained graph model combining
masked graph autoencoding, temporal relational diffusion over the
service-dependency graph, and causal attention over multi-hop paths to
separate likely causes from correlated downstream effects. This is the
paper closest in domain and in intent to distinguish root cause from
symptom under propagation -- but it is still a model of \emph{which
services call which and how anomaly propagates structurally}, learned from
graph topology and telemetry. It does not represent liquidity-lock
reservation state, AML-hold state, idempotency-key collision state, or
settlement-finality state as inputs at all. Section~\ref{sec:eval}'s
central empirical result is that this specific gap is exploitable: on
incidents deliberately constructed so a downstream service's telemetry
anomaly is louder than the true upstream root's own signal, a
topology/diffusion-style structural model has exactly the failure mode
ServiceGraph-FM's own causal-attention mechanism is designed to resist --
and payment-state evidence, which no structural model in this line of work
consumes, is what actually resolves it.

\subsection{Synthetic Payment and AML Transaction Simulation}
PaySim~\cite{paysim2016} derives payment amount and timing distributions
from real aggregated mobile-money statistics and injects labeled fraud
without using real entity identities. IBM AMLworld~\cite{amlworld2023} is
agent-based and typology-driven (fan-in/fan-out, structuring, layering)
with an explicit anonymization and ethical-use stance. The generator in
Section~\ref{sec:data} is grounded methodologically in this line of work
and, notably, was not originally built that way: an earlier internal
version used real corporate names and real sanctioned-individual names
drawn from public sanctions lists for realism. Identifying and correcting
this during development, rather than after, is stated here directly as a
design correction rather than omitted -- it strengthens the methodology
section, not weakens it, and is the same anonymization discipline
AMLworld's ethics stance establishes as the norm for this kind of dataset.

\subsection{Payment System Operational Resilience}
The Bank for International Settlements' Core Principles for Systemically
Important Payment Systems~\cite{bis2001cpss} and the IMF's more recent
survey of operational resilience in digital
payments~\cite{imf2021resilience} establish that a disruption at a single
participant or processing stage can have consequences well beyond that
stage -- the policy-level statement of the same structural fact this
paper's fault taxonomy exercises algorithmically. This body of work is positioned deliberately as policy and
continuity-level, not competing with the algorithmic RCA literature
surveyed above: it motivates \emph{why} diagnosing payment-pipeline
failures quickly matters at a systemic level, while this paper addresses
\emph{how} to diagnose them correctly once a failure has occurred.

\textbf{Note on a retracted citation:} an earlier literature pass surfaced a
candidate paper (``Han, latency-aware payment RCA via heterogeneous graph
attention,'' \emph{Computer Science Bulletin}) that was carried as
tentative prior art. A direct web search (2026-08-27) for this title,
venue, and author combination returns no matching paper, indexed venue, or
author/institution trail -- it does not check out and is not cited. Its
role as ``closest prior art'' is filled instead by
ServiceGraph-FM~\cite{servicegraphfm2026} above, which is independently
verified (resolves via DOI on MDPI).

\section{System Architecture}
\label{sec:system}

ClearFlow is a five-service event-driven payment pipeline built on Spring
Boot 3.3.2 / Java 21, coordinated via Apache Kafka:
\texttt{gateway} $\rightarrow$ \texttt{validation-enrichment} $\rightarrow$
\texttt{aml-compliance} $\rightarrow$ \texttt{routing-execution}
$\rightarrow$ \texttt{settlement}, with \texttt{fraud-scoring} and
\texttt{audit} as parallel, non-blocking consumers of the initiation event.
Payments are submitted as ISO~20022 pacs.008-shaped requests
(\texttt{instructionId}, \texttt{endToEndId}, \texttt{uetr}) and processed
through 18 Kafka topics with per-stage idempotency (Redis-backed) and an
outbox relay at the gateway.

\texttt{fraud-scoring} is a real LightGBM classifier (11 engineered
features: amount, temporal, country/currency risk, and real-time
transaction-velocity signals) served over HTTP with a circuit-breaker
fallback to a heuristic scorer, not a synthetic stub -- but it is
explicitly not a research contribution of this paper. It exists only to
produce realistic fraud-adjacent payment state (\texttt{fraudScore},
\texttt{riskBand}) as a precondition for fault scenarios, keeping this
paper's claim scoped to payment-state-aware RCA rather than fraud
detection itself. For completeness: the deployed model is trained on a
public mobile-money fraud dataset (PaySim) with genuine, extreme class
imbalance (186 fraud examples in 400,000 training rows) and reaches AUC
0.62 on a held-out test set -- modest but honest for this rarity. A
retrain attempt on a richer dataset (more positive examples, real
transaction-velocity signal) reached AUC 0.86 by that metric alone, but
was found, before deployment, to produce a badly miscalibrated,
near-bimodal score distribution on realistic inputs (roughly half of
simulated ordinary payments scored above 0.99), and was rolled back
rather than shipped -- reported here as a concrete illustration that a
single ranking metric does not certify a model fit for a live
decisioning threshold.

Three transactional-state mechanisms are central to this paper's RCA
contribution:
\begin{itemize}
  \item \textbf{Liquidity reservation} -- \texttt{routing-execution} takes a
  pessimistic lock (\texttt{SELECT FOR UPDATE}) against a nostro account
  balance for the duration of routing and settlement.
  \item \textbf{AML holds} -- \texttt{aml-compliance} screens against a
  sanctions/PEP list and can block a payment from progressing, independent
  of downstream service health.
  \item \textbf{Idempotency and settlement finality} -- retries are
  deduplicated by idempotency key; once \texttt{SETTLEMENT\_COMPLETE} is
  emitted, the transaction is final and cannot be rolled back, only
  compensated.
\end{itemize}

\todo{Add architecture diagram (Fig.~1): service graph + Kafka topics +
state-holding points. Verified live end-to-end 2026-08-17: single payment
traced gateway$\to$validated$\to$AML-cleared$\to$SEPA\_INSTANT
routed$\to$settled, sub-second, correlation-ID traceable across all 5
services -- use this trace as the walkthrough example in the figure
caption.}

\section{Synthetic Payment and Fault Data Generation}
\label{sec:data}

A synthetic corpus generator produces the fast-iteration dataset used
alongside the live evaluation of Section~\ref{sec:eval}. An earlier version
of this generator used real corporate and sanctioned-individual names for
realism; this was identified as a liability during development and
replaced entirely with fully synthetic entity identifiers, following the
anonymization precedent set by IBM AMLworld~\cite{amlworld2023} -- no real
company, individual, or sanctioned-entity name appears anywhere in the
generator or its output.

\textbf{Scale.} 11{,}000 persistent debtor/creditor account profiles
(country, home currency, risk tier, PEP flag, behavioral baseline) support
50{,}000 generated payments and a derived causal event timeline of
approximately 438{,}000 state-transition events (8.76 events/payment on
average). Rail selection is conditioned on currency (SEPA rails carry only
EUR, FEDWIRE/FEDACH/CHIPS only USD, etc.) and each rail enforces a real
amount ceiling at selection time, following PaySim-style~\cite{paysim2016}
log-normal amount distributions with type-specific medians. AML scenarios
use AMLworld-style typology labels (structuring, layering, integration);
an approximately 1\% embargo-corridor injection rate exercises the same
compliance logic the live system's AML-hold gate implements
(Section~\ref{sec:system}), and idempotency-collision rows are generated by
cloning an existing payment's debtor/amount/creditor fields and
recomputing the same SHA-256 idempotency key the live gateway derives
(Section~\ref{sec:system}), so the synthetic dataset's idempotency
semantics match the live system's exactly rather than approximating them.

\textbf{Incident corpus.} On top of the payment corpus, controlled
incidents are injected across the same four fault families used in the
live evaluation (infra, payment-domain, cross-domain, confounded),
balanced across fault types rather than left to occur at whatever rate
they would naturally, with hierarchical ground truth (root service,
component, causal event, propagation path/depth, severity) and
temporal-difficulty labels. Confounded incidents are constructed so a
downstream symptom service's metric deviation is quantitatively larger
than the true root's own -- verified directly against the generated data,
not merely asserted by design, for every confounded incident in the
corpus. Ground-truth incident metadata and the incident-to-payment linkage
used for scoring are written to files kept separate from the evidence
surface any RCA method reads, so a method cannot trivially read its own
answer key; a dedicated graph-construction check verifies programmatically
that no evidence tier a method is restricted to ever exposes a ground-truth
node or edge, rather than relying on the separation being followed by
convention.

\section{Fault Injection and RCA Ground Truth}
\label{sec:faults}

An earlier version of this evaluation used a synthetic corpus generator
producing labeled ground truth by construction (payment/account
distributions grounded in PaySim~\cite{paysim2016} and
AMLworld~\cite{amlworld2023}, fault injection with hand-set propagation
paths and severities). That corpus remains useful for fast iteration and is
reported for comparison in Section~\ref{sec:eval}, but its ground truth is
generated, not measured -- an evaluation on it alone cannot distinguish "the
method works" from "the method recovers exactly the assumptions the
generator encoded." This section describes the live fault-injection
harness built to close that gap: it triggers \emph{real} faults against the
running system of Section~\ref{sec:system} and derives ground truth from
what actually happens, not from a specification.

\subsection{Fault taxonomy and live-triggerability}
Ten fault types across the four families standard to this evaluation
protocol (infra, payment-domain, cross-domain, confounded) are triggered
through two real mechanisms:

\begin{itemize}
  \item \textbf{Process crash} (infra, cross-domain, confounded families):
  an admin endpoint kills and cold-restarts the target service by port
  (\texttt{lsof|kill -9}, relaunch the jar), producing a genuine outage
  under concurrent live traffic. Confounded incidents pair a crash on the
  true root service with sustained traffic through a service positioned to
  show backpressure as a louder, correlated symptom -- e.g.\ a
  \texttt{validation-enrichment} crash upstream of \texttt{gateway}, which
  absorbs the resulting queueing/retry pressure and can show a larger raw
  telemetry deviation than the service that actually failed.
  \item \textbf{Payment-domain triggers}: an AML hold is triggered by
  submitting a payment against a real sanctions-list entity (confirmed via
  the live compliance-review endpoint, not assumed); an idempotency
  collision is triggered by resubmitting an accepted payment's exact
  payload within the gateway's deduplication window (confirmed via the
  real HTTP 409 response).
\end{itemize}

Two fault types specified in an earlier design pass were found, empirically,
not to be live-triggerable through the running system as built, and are
excluded from the live corpus rather than simulated: a settlement-finality
violation is intercepted by an idempotency guard and a database existence
check before the finality logic it is meant to test ever executes, and a
liquidity-lock-stuck fault requires a hand-constructed Kafka-level message
for which no REST path currently exists. Reporting this honestly is itself
part of the contribution: a fault taxonomy validated only on paper, without
checking whether each fault type is actually reachable through the system
under test, risks silently testing the generator's assumptions rather than
the system's real behavior.

\subsection{What ground truth is known vs.\ measured}
Because the harness chooses which fault to trigger and when, \texttt{fault\_type},
\texttt{root\_service}, and \texttt{injection\_time} are known with certainty --
this is the one property live triggering shares with a synthetic generator.
Everything else that would be asserted by a generator is instead measured
post-hoc from real evidence: \texttt{n\_affected\_payments} is counted from
real payment submissions inside the incident window;
\texttt{propagation\_depth} is measured as the number of services (beyond
the root) whose real error-rate z-score exceeds a 2$\sigma$ deviation from
their own pre-incident baseline during the window; \texttt{severity} is
derived from the observed affected-payment count, not chosen in advance.
This measurement is itself imperfect and its limitations are stated
directly in Section~\ref{sec:discussion} rather than left implicit.

A methodological hazard specific to running many fault injections against
one live system in sequence is worth naming explicitly, because it produced
a real bug during this work: the pre-incident baseline window used to
compute z-scores must not overlap another incident's own effects, or every
"quiet baseline" is contaminated by nearby injected faults. Incidents are
therefore spaced with an enforced cooldown exceeding the baseline lookback
window, and this spacing is logged and checked, not assumed.

\section{RCA Method}
\label{sec:method}

Three methods are compared, each explicitly restricted to a named evidence
tier so that an accuracy difference can be attributed to what a method is
allowed to see, not to one method simply having more information. None of
the three ever reads ground-truth incident labels or the incident-to-payment
linkage held out for scoring.

\subsection{Baseline: telemetry and topology}
\texttt{loudest\_metric\_baseline} ranks services purely by real-time
error-rate anomaly: for each service, a z-score of its error rate during the
incident window against its own pre-incident baseline (evidence tier G2 --
telemetry only, no topology, no payment state). This is the "find the
biggest spike" baseline every RCA benchmark in Section~\ref{sec:related}
implicitly competes against, and the one a confounded incident is
specifically constructed to defeat.

\texttt{graph\_topology\_baseline} (tier G0--G2) adds pipeline topology as a
tie-breaker: when two services' z-scores fall within a bounded margin of
each other, the more upstream service in the fixed pipeline order is
preferred. Topology is deliberately restricted to breaking close calls, not
overriding a clear telemetry signal -- an earlier, more aggressive version
that let topology override z-score rank unconditionally was found to score
\emph{worse} than telemetry alone on confounded incidents (a hard
upstream-preference rule is wrong whenever backpressure makes a downstream
symptom louder, which is exactly what a confounded incident constructs),
and is documented here as a negative result rather than discarded silently.

\subsection{Payment-state-aware RCA}
\texttt{payment\_aware\_rca} (tier G0--G3) adds the transactional-state
mechanisms of Section~\ref{sec:system} as evidence: for payments active
during the incident window, the fraction showing an AML hold, a stuck
liquidity reservation, a detected idempotency collision, a failed
settlement, or (a generalization applying to every crash-based fault, since
a process crash touches no payment's own state fields) having stalled
before reaching a given pipeline stage's completion event entirely. A
domain-state fraction elevated above a fixed threshold is treated as
\emph{decisive} -- ranked ahead of telemetry -- rather than as an additive
bonus on top of the z-score. This distinction matters mechanically: a fixed
additive term cannot reliably outrank a z-score gap that scales with
incident severity, which is precisely what a confounded incident's louder
downstream symptom does at high severity; only a decisive override is
structurally capable of correcting it. With no elevated payment-state
signal, the method falls back to \texttt{graph\_topology\_baseline}'s
ranking unchanged.

Two implementation-level findings, surfaced by tracing specific
misclassified incidents against real (not synthetic) evidence rather than
by tuning against an aggregate accuracy number, are reported here because
they generalize beyond this specific system: first, a state-flag defined
by a single evidence snapshot (payment reserved and not yet settled) is not
specific evidence of a stuck fault when a system has genuine settlement
latency -- it is also the normal appearance of any payment that simply has
not finished yet, and must be gated by how long the payment has actually
been observed in that state, not by its state at one instant. Second, this
generalizes directly to the "never reached a stage at all" case: a crash-based
fault produces no state anomaly whatsoever in the schema this paper's
system tracks, and the only signal available is whether a payment's
expected completion event for a given stage occurred within a reasonable
multiple of that stage's normal latency. Both findings were verified not to
change the earlier, generator-based evaluation's results at all (Section~\ref{sec:eval}),
which is itself evidence they correct a real gap between synthetic and live
evidence rather than fitting noise in the live sample they were derived
from.

\subsection{LLM baseline}
A fourth method, evaluated at the same G0--G2 evidence tier as
\texttt{graph\_topology\_baseline}, tests whether a general-purpose language
model reasoning over the same topology and telemetry text outperforms the
formulaic heuristic. The model is given the pipeline order, each service's
z-score, and an explicit description of the backpressure/loud-symptom
failure mode confounded incidents are constructed to test, then asked to
rank all five services. This is a controlled comparison of reasoning
strategy given identical evidence, not yet a claim about whether an LLM
given richer (G0--G3) evidence would do better -- that comparison is
addressed directly by the next two methods, which isolate the evidence
axis and the reasoning axis from each other rather than changing both at
once.

\subsection{Static LLM with G0--G3 evidence}
A fifth method holds reasoning mode fixed (the identical static, one-shot
prompt structure as the G0--G2 baseline) and changes only the evidence
tier: the prompt additionally includes the same six payment-domain state
fractions Section~\ref{sec:method}'s formula computes (AML hold, liquidity
dwell, idempotency, settlement failure, validation retry/stall), presented
as text, with no tool access. Compared directly against the G0--G2 LLM
baseline, this isolates whether richer evidence alone -- independent of
how it is used -- improves LLM-based diagnosis.

\subsection{Agentic RCA}
A sixth method addresses the reasoning axis directly, with real
tool access rather than a richer static prompt: given the same G0--G2
topology/telemetry context as the LLM baseline, the model is also given
two tools backed by the live MCP gateway's real per-payment endpoints --
retrieve a payment's full stage-by-stage timeline, or its AML screening
detail -- for any of the payments active in the incident window, and may
call them iteratively (up to 4 rounds) before answering. This is a
materially different agent design than a static prompt: the model decides
for itself which of the window's payments are worth inspecting and what
to look at, rather than being handed a precomputed summary. MCP's
aggregate endpoints (systemic health, active alerts) are deliberately
excluded as tools: they aggregate over a window counted back from
wall-clock query time, not from the incident's own historical timestamp,
and would mix in unrelated later activity for incidents queried hours
after they occurred -- only the per-payment endpoints are historically
exact regardless of when the agent runs.

\section{Evaluation}
\label{sec:eval}

\subsection{Protocol}
AC@1/AC@3 are reported per the RCAEval convention~\cite{rcaeval2024}, with
one deliberate departure justified in Section~\ref{sec:method}: AC@5 is
omitted, since this system has exactly 5 candidate services and AC@5 is
therefore 1.0 for every method by construction, not an evaluation result.
Every AC@k number is reported with a Wilson 95\% confidence interval, and
family-level and depth-level strata are reported with their $n$ directly
alongside the mean rather than as a bare fraction. Methods are compared
via exact McNemar paired significance testing on per-incident AC@1
hit/miss; a comparison is reported only when at least 10 discordant pairs
exist, and is explicitly marked as insufficient otherwise rather than
implying a precision the sample does not support.

The headline result uses a genuine held-out split: the 17 live incidents
collected while the dwell-time gate and stage-stall generalization
(Section~\ref{sec:method}) were being developed and diagnosed against
specific misclassified cases form a frozen \emph{dev set}; only the 25
incidents collected \emph{after} both fixes were finalized in code are
scored for the comparison below. Two reference-floor methods are scored
alongside the four real methods: \texttt{random\_baseline} (uniform random
ranking, AC@1 expectation 0.2) and \texttt{majority\_baseline} (always
guess the most common root service, fit only on the disjoint dev set).

\begin{table}[h]
\centering
\caption{Root Cause Identification Accuracy, held-out live incidents
($n=25$, dev-set-derived, never used to develop the method)}
\begin{tabular}{lcccc}
\toprule
Method & infra & pay.-dom. & cross-dom. & \textbf{AC@1} \\
 & (n=12) & (n=6) & (n=4) & (n=25) \\
\midrule
random\_baseline           & 0.17 & 0.17 & 0.25 & 0.16 \\
majority\_baseline         & 0.25 & 0.50 & 0.50 & 0.32 \\
llm\_rca\_baseline (G0--G2, static) & 0.25 & 0.50 & 0.25 & 0.28 \\
llm\_rca\_g3\_baseline (G0--G3, static) & 0.25 & 0.67 & 0.00 & 0.28 \\
loudest\_metric (G2)        & 0.58 & 0.33 & 0.50 & 0.44 \\
graph\_topology (G0--G2)    & 0.58 & 0.50 & 0.50 & 0.48 \\
agentic (G0--G3, LLM+MCP tools) & 0.58--0.67 & 0.50 & 0.50 & 0.44--0.52 \\
\textbf{payment-aware (G0--G3, ours)} & \textbf{0.67} & \textbf{0.83} & 0.50 & \textbf{0.60} \\
\bottomrule
\end{tabular}
\end{table}

Confounded ($n=3$) is omitted from the table above and discussed
separately: every method, including \texttt{random\_baseline}, scores 0.0
on the held-out confounded incidents -- reported in full in
Section~\ref{sec:eval-confounded}.

\subsection{Held-out result}
\texttt{payment\_aware\_rca} reaches overall AC@1 0.60 (95\% CI
0.41--0.77) on the held-out set, ahead of \texttt{graph\_topology\_baseline}'s
0.48 (CI 0.30--0.67) and \texttt{loudest\_metric\_baseline}'s 0.44 (CI
0.27--0.63). McNemar paired testing against \texttt{random\_baseline}
reaches significance ($p=0.0034$, 13 discordant pairs) -- the only
comparison in this evaluation with enough discordant pairs to clear the
10-pair reporting threshold set in Section~\ref{sec:eval}'s protocol. The
comparisons against \texttt{loudest\_metric\_baseline} (6 discordant
pairs), \texttt{graph\_topology\_baseline} (5), and \texttt{llm\_rca\_baseline}
(8) all fall below that threshold; the point estimates favor
\texttt{payment\_aware\_rca} in every case, but no significance claim is
made for any of them. \texttt{payment\_domain} is the strongest and most
consistent result (0.83 vs.\ 0.50 and 0.33) -- the family the payment-state
signals were specifically built to disambiguate, and the family where the
observed effect size is least likely to be an artifact of the small held-out
sample.

The LLM baseline result from an earlier (dev-set) pass replicates on
genuinely fresh data: 0.28 overall, below \emph{both} non-LLM baselines
given identical G0--G2 evidence. This is the one dev-set finding that
held up unchanged on held-out incidents, and is accordingly the result
this paper trusts most among the four methods' comparisons.

\textbf{Evidence vs.\ reasoning: a clean, direct answer.}
\texttt{llm\_rca\_g3\_baseline} -- the identical static prompting strategy
as the G0--G2 baseline, given the same G0--G3 payment-state fractions the
formula uses -- scores \emph{identically} overall: 0.28 vs.\ 0.28. Handing
a static LLM the exact same richer evidence a formula uses to reach 0.60
produces \emph{zero net change}. This is not uniformly null underneath,
which is itself informative: by family, the added evidence helps
\texttt{payment\_domain} (0.50 $\to$ 0.667, exactly the family those
fractions are literally about) and actively hurts \texttt{cross\_domain}
(0.25 $\to$ 0.00, where the added dense numeric text appears to distract
more than inform), and these effects cancel to an identical pooled result
at $n=25$ -- reporting the family breakdown rather than only the pooled
number is what surfaces this instead of averaging it away. Tool access,
by contrast, improves broadly and specifically on \texttt{infra} (0.25 $\to$
0.58--0.67): given the same information but the reasoning mode to actively
decide what to inspect rather than passively receive it, the agent
recovers most of the gap to the formulaic method. Together these results
answer the question directly: \texttt{payment\_aware\_rca}'s advantage over
the LLM-based methods is not explained by access to richer evidence -- the
static LLM had identical evidence and did not improve -- it is explained
by the formula's decisive, structured use of that evidence
(Section~\ref{sec:ablation}), which neither static prompting nor
currently-scoped agentic tool-use fully replicates.

\textbf{Tool access closes most of the gap to the formulaic method.} Given
the same G0--G2 context as the static LLM baseline plus real tool access
to the live system, the agentic method reaches 0.44--0.52 overall AC@1
(two repeated runs on the identical held-out set differ by 0.08 despite
temperature 0 -- unlike every formulaic method, which is bit-for-bit
reproducible on the same input, the agentic method's real-time tool-call
round-trips introduce genuine run-to-run variance, itself worth stating
rather than averaging away). That is more than double the static LLM
baseline's 0.28 and matches or exceeds both non-payment-state baselines,
purely from letting the model decide what real evidence to look at rather
than reasoning over a fixed telemetry summary -- while still trailing
\texttt{payment\_aware\_rca}'s 0.60. On \texttt{infra} specifically the
agent ties the payment-aware method exactly (0.67 vs.\ 0.67).

\textbf{Dollar-exposure weighting surfaces a real gap, confounded by
fault duration.} Weighting each held-out incident's AC@1 hit/miss by the
real dollar value of payments active in its window (extracted from live
Elasticsearch logs, illustrative FX to USD) drops
\texttt{payment\_aware\_rca}'s exposure-weighted AC@1 to 0.33 -- well
below its unweighted 0.60, and the largest such gap of any method tested.
Investigated before reporting it as ``fails on large transactions'':
sorting held-out incidents by exposure shows a clean bimodal split
(\$900K--\$5.5M for every crash-mechanism incident vs.\ \$4K--\$47K for
every payment-domain-trigger incident), which traces directly to
injection duration (20--30s crash faults accumulate far more concurrent
payments than the 5s payment-domain triggers) rather than to transaction
size specifically -- and correlates with \texttt{confounded}/\texttt{cross\_domain}
already being the hardest families independent of dollar value. The
honest reading: there is a real accuracy gap by fault family that
dollar-weighting surfaces starkly, but confirming it is specifically about
transaction size (not duration) needs exposure normalized per second of
window rather than summed over it -- not built here, and stated as an
open question rather than an unqualified finding either way.

\subsection{The confounded result does not replicate}
\label{sec:eval-confounded}
An earlier analysis pass on the dev set found \texttt{payment\_aware\_rca}
recovering strongly on confounded incidents (0.667 AC@1, beating both
baselines) after the stage-stall generalization was added. On the 3
confounded incidents in the held-out set, every method -- including
\texttt{random\_baseline} -- scores 0.0. This is reported exactly as
observed, not smoothed over: it is precisely the failure mode a held-out
split exists to catch, and it means the confounded-family improvement
observed during development does not demonstrably generalize. Whether
this reflects a real, family-specific limitation of the current signal set
or is simply too small an $n$ (3 incidents) to conclude anything is left
genuinely undetermined rather than resolved either way -- a larger
confounded-specific live sample is the direct next step this result
implies.

\subsection{Ablation: which mechanism actually drives the result}
\label{sec:ablation}
\texttt{payment\_aware\_rca} combines two mechanisms (Section~\ref{sec:method}):
6 payment-domain-specific state fractions (AML hold, liquidity dwell,
idempotency, settlement failure, validation retry/stall) and a generalized
5-stage completion-stall signal added specifically to recover accuracy on
crash-based faults. Three ablations isolate each mechanism's actual
contribution on the same 25 held-out incidents, rather than reporting the
combined method as a black box:

\begin{table}[h]
\centering
\caption{Ablation, held-out AC@1 ($n=25$)}
\begin{tabular}{lc}
\toprule
Variant & AC@1 \\
\midrule
fracs-only (stall signal disabled) & 0.60 \\
stall-only (all 6 fracs disabled) & 0.40 \\
no dwell-time gate (pre-fix behavior) & 0.56 \\
\textbf{full method} & \textbf{0.60} \\
\bottomrule
\end{tabular}
\end{table}

The result is more specific, and more useful, than "the method works":
\textbf{fracs-only exactly matches the full method (0.60 = 0.60)} -- the
generalized stall signal, despite producing the dev-set's headline
confounded recovery (Section~\ref{sec:eval-confounded}), contributes
nothing measurable beyond the payment-domain fractions on genuinely
unseen incidents. \textbf{stall-only alone underperforms even
\texttt{graph\_topology\_baseline}} (0.40 vs.\ 0.48), meaning the
generalization is not simply weaker than the fracs -- on its own it is
actively worse than not using payment-state evidence at all. This
directly explains, rather than merely correlates with, why confounded did
not replicate: the specific mechanism responsible for that dev-set result
is the one ablation shows does not generalize. The dwell-time gate ablation
confirms that fix is real but more modest on this held-out set than its
dev-set diagnosis suggested (0.56 vs.\ 0.60) -- expected, since the
incident it was traced from was itself in the dev set, not held out.

The paper's validated, mechanism-isolated claim is therefore narrower and
more defensible than the headline number alone suggests: payment-domain
transactional state (Section~\ref{sec:system}'s liquidity/AML/idempotency
mechanisms), not process-crash-derived stage-completion evidence, is what
this evaluation actually supports as a causal RCA signal.

\section{Experiment 2: State-Conditioned Failure Propagation}
\label{sec:exp2}

Experiment 1 (Sections~\ref{sec:eval}) asks a diagnostic question: given a
cascade, which service caused it. This section asks a different,
interventional question that only a live system with real controllable
faults can pose at all: does the \emph{same} infrastructure fault produce
a \emph{different} cascade depending on what real, persistent
payment-domain state already exists in the system when it hits. Every
prior-art system in Section~\ref{sec:related}, ServiceGraph-FM included,
reasons over passive, already-collected telemetry -- none can run this
comparison, because none control the intervention.

\subsection{Design}
The same fault (\texttt{DB\_TIMEOUT} on \texttt{settlement}) is triggered
under two conditions. \textbf{Condition A (clean background)} reuses the 5
\texttt{DB\_TIMEOUT} incidents already collected for Experiment 1 as-is --
not re-triggered, since Experiment 1's own data already is this fault
type's clean-background condition. \textbf{Condition B (AML-held
background)}: 2 real AML holds are triggered first and deliberately left
unresolved -- genuinely persistent live state (a held payment stays
\texttt{HOLD} until the compliance resolve endpoint is explicitly called,
which this experiment never does), not a synthetic label -- then the same
\texttt{DB\_TIMEOUT} fault is triggered while those holds remain active,
repeated for 4 incidents. Both conditions are measured on the identical
pipeline used throughout Experiment 1 (real error-rate z-scores,
\texttt{measure\_propagation\_depth}, real per-payment amount/currency for
dollar exposure) so no separate, ad-hoc metric definitions are introduced
for this comparison.

\subsection{Result: an honest null finding}
\begin{table}[h]
\centering
\caption{Experiment 2: Condition A vs.\ B, raw values (n too small for significance testing)}
\begin{tabular}{lcc}
\toprule
Metric & Condition A ($n=5$) & Condition B ($n=4$) \\
\midrule
n\_affected\_payments (mean) & 15.0 & 14.5 \\
propagation\_depth (mean) & 1.00 & 1.00 \\
exposure\_usd (mean) & \$3.75M & \$4.32M \\
\bottomrule
\end{tabular}
\end{table}

At this sample size, a persistent AML-hold background does not detectably
change how a settlement database failure cascades.
\texttt{propagation\_depth} is identically 1 across every single incident
in both conditions. \texttt{n\_affected\_payments} ranges overlap almost
completely (11--19 vs.\ 13--16). \texttt{exposure\_usd} ranges overlap
heavily despite Condition B's mean running somewhat higher; at $n=4$--$5$
this is not distinguishable from noise, and no significance test is
attempted -- fabricating one this sample cannot support would undercut the
methodological discipline the rest of this paper is built on.

Two readings are stated, and neither is asserted as the answer. This
could be a real, operationally reassuring finding: the system's fault
isolation held, and an unrelated compliance hold elsewhere did not
amplify an unrelated infrastructure failure's blast radius. Or $n=4$--$5$
per condition may simply be too small to detect a real effect, compounded
by \texttt{propagation\_depth}'s own measurement ceiling -- this live
system's crash-recovery is fast enough that depth rarely exceeds 1 for
any infra fault under either condition (Section~\ref{sec:eval}), limiting
how much room this specific metric has to show a state-conditioning
effect even if one exists. A larger Condition-B sample and a fault type
whose depth already varies more under clean conditions are the concrete
next steps this null result implies. Reporting a carefully-checked null
result from an experiment no passive-dataset methodology in this
literature could even run is more useful to the field than a convenient
positive result would have been.

\section{Discussion}
\label{sec:discussion}

\subsection{Threats to validity}
\label{sec:threats}

\textbf{Sample size.} The live evaluation in Section~\ref{sec:eval} runs at
a small number of incidents per fault family. Every result is reported with
a Wilson 95\% confidence interval specifically so this is not obscured by a
bare mean; a family-level accuracy difference of one or two incidents is
not distinguishable from chance at this $n$, and is reported as such rather
than as a finding. A larger live corpus, not a synthetic one, is the direct
next step this limitation implies, since the entire motivation for building
the live harness (Section~\ref{sec:faults}) was that synthetic ground truth
cannot answer whether a method's synthetic-benchmark result transfers.

\textbf{Tautological ground truth for crash-based faults.} For every
infra/cross-domain/confounded incident, \texttt{root\_service} is, by
construction, the process the harness chose to kill. This is real
infrastructure and real telemetry noise, but it is \emph{not} genuine
diagnostic ambiguity of the kind a real operator faces when the root cause
is actually unknown -- the harness always knows the answer it is testing
against. This is the same limitation present in essentially all fault-
injection-based RCA benchmarks, including the ones surveyed in
Section~\ref{sec:related}, and is stated plainly here rather than allowed
to be implied by the word ``real.''

\textbf{Threshold selection.} The payment-state elevation threshold and the
dwell-time gates in Section~\ref{sec:method} are domain-motivated constants
justified by measured normal-operation latencies (e.g.\ observed real
settlement completion in milliseconds), not values fit or cross-validated
against held-out data. A genuine train/validation split for these specific
constants, distinct from the incident-level dev/held-out split already used
for the headline comparison, is future work once the live corpus is large
enough to support it without starving the held-out test set.

\textbf{Single-organization topology.} The pipeline models one
organization's internal payment processing, not a multi-participant/
consortium settlement network where root cause can originate outside the
observed system entirely. The payment-state mechanisms evaluated here
(liquidity locks, AML holds, idempotency, finality) are internal to a single
processor's pipeline; whether they generalize to a network-level topology
is untested.

\textbf{Measured, not asserted, propagation depth.} \texttt{propagation\_depth}
is itself a z-score-threshold heuristic computed from the same telemetry
the RCA methods consume, not an independently established ground truth.
Depth-stratified results in Section~\ref{sec:eval} should be read as
descriptive of what the same measurement pipeline reports, not as evidence
of a causally verified depth effect.

\subsection{Why report this at all}
\label{sec:whyreport}
Every limitation above was found by directly auditing this paper's own
evaluation methodology after an initial round of results looked cleaner
than the sample size could support -- an earlier draft of
Section~\ref{sec:eval}'s live result reported a family-level accuracy
comparison derived from incidents whose baseline windows were, it turned
out, contaminated by adjacent fault injections, and a fix validated only
against the same incidents used to diagnose it. Both were caught and
corrected before this version, not smoothed over, and the corrected
evaluation pipeline (Wilson CIs, a genuine held-out split, reference floor
baselines, paired significance testing) is the pipeline
Section~\ref{sec:eval} actually reports. This is stated explicitly because
a systems paper whose own author found and disclosed these gaps is a
stronger, not weaker, paper than one where a reviewer finds them first.

\section{Conclusion}
\label{sec:conclusion}

Payment-specific transactional state -- liquidity locks, AML holds, and
idempotency collisions -- is a causal factor in how failures cascade
through a payment pipeline, and existing microservice RCA methods, general
or LLM-based, do not model it. On real incidents triggered against a live,
running system and scored against a held-out set disjoint from the
incidents used to develop the method, a payment-state-aware RCA method
reaches AC@1 0.60, ahead of both a topology-aware and a telemetry-only
baseline, and significantly ahead of random guessing. An ablation, not
just the headline number, is what makes this claim defensible: the
payment-domain state fractions alone reproduce the full method's result
exactly, while a generalized process-crash completion-stall signal
contributes nothing beyond them and underperforms alone -- the paper's
validated contribution is specifically payment-domain transactional
state, not a general theory of crash-fault evidence, and the confounded
family's dev-set improvement not replicating on held-out data is the
direct, explained consequence of that same finding rather than an
unrelated gap. One implementation lesson generalizes beyond this system
regardless of that scope: a single-snapshot state flag is not specific
evidence of a stuck fault without a dwell-time gate, verified to matter on
live data (0.56 vs.\ 0.60 without it) while leaving the synthetic
benchmark's numbers completely unchanged -- evidence it corrects a real
synthetic-to-live gap rather than fitting noise in the sample it was
derived from. Extending payment-domain-specific state evidence to the
fault families it does not yet cover, rather than generalizing the crash
signal further, and extending the LLM baseline to the same G0--G3 evidence
tier the payment-aware method uses, are the direct next steps this
evaluation's own ablation points to.

A separate, interventional experiment (Section~\ref{sec:exp2}) --
possible only because the evaluation system is live and controllable, not
a passive dataset -- found no detectable difference in cascade
propagation between two real, controlled payment-state backgrounds for
the fault type tested. That null result is reported with its own
limitations stated directly, not smoothed over: a larger sample and a
fault type with more room to show a state-conditioning effect are the
concrete next steps it implies, and the experiment's own value is that no
prior-art system in this literature could have run it at all.

\bibliographystyle{IEEEtran}
\bibliography{references}

\end{document}
